%% TREŚĆ ZADANIA, NUMER, IMIĘ I NAZWISKO

\begin{problem}{2.23}{Łukasz Bożyk}
Liczby dodatnie $a$, $b$, $c$ spełniają nierówność \begin{equation}\label{2.23:teza}
% etykiety równań proszę zaczynać od numeru zadania jak zilustrowano wyżej, żeby po wrzuceniu wszystkich prac do jednego pliku nie było kolizji
\frac{a^2+b^2+c^2}{ab+bc+ca}<2.\end{equation}
Wykaż, że istnieje trójkąt o bokach $\sqrt{a}$, $\sqrt{b}$, $\sqrt{c}$.
\end{problem}

%% MIEJSCE NA ZAWARTOŚĆ OPRACOWANIA

\paragraph{Rozwiązanie 1.} Załóżmy bez straty ogólności, że $a\leq b\leq c$. Teza zadania przybiera wówczas postać $\sqrt{a}+\sqrt{b}>\sqrt{c}$. Skoro obie strony są dodatnie, to równoważnie mamy
\begin{align*}
    a+b+2\sqrt{ab}&>c,\\
    2\sqrt{ab}&>c-a-b.
\end{align*}
Jeżeli $c-a-b<0$, to ostatnia nierówność jest oczywiście spełniona (w~tym wypadku teza zadania jest więc spełniona). Przyjmując dodatkowe założenie $c-a-b\geq 0$, możemy ponownie podnieść nierówność stronami do kwadratu, uzyskując \emph{równoważną} postać tezy (w~tym przypadku):
\begin{align*}
    4ab&>c^2+a^2+b^2+2ab-2ac-2bc,\\
    2(ab+bc+ca)&>a^2+b^2+c^2,\\
    2&>\dfrac{a^2+b^2+c^2}{ab+bc+ca}.
\end{align*}
Otrzymaliśmy nierówność \eqref{2.23:teza}. Ponieważ przekształcenia były równoważne, więc rozwiązanie jest zakończone.

\paragraph{Rozwiązanie 2.}
    Przyjmijmy oznaczenia $x=\sqrt{a}$, $y=\sqrt{b}$, $z=\sqrt{c}$. Daną nierówność można przepisać jako
    \[2x^2y^2+2y^2z^2+2z^2x^2-x^4-y^4-z^4>0,\]
    czyli równoważnie
    \begin{equation}\label{2.23:iloczyn}(x+y+z)(x+y-z)(y+z-x)(z+x-y)> 0.\end{equation}
    Czynnik $x+y+z$ jest dodatni, więc spośród pozostałych trzech czynników albo wszystkie są dodatnie (i~zadanie jest rozwiązane, gdyż to dokładnie oznacza, że istnieje trójkąt o~bokach $x$, $y$, $z$), albo dwa są ujemne. Jednak drugi przypadek nie może mieć miejsca: gdyby np. $x+y-z<0$ oraz $y+z-x<0$, to dodając te dwie nierówności stronami, uzyskalibyśmy $2y<0$ --- sprzeczność.

\paragraph{Rozwiązanie 3.} Korzystając (dwukrotnie) z~faktu, że nierówność $|x-y|<z$ jest równoważna układowi nierówności $-z<x-y<z$, przekształcamy daną nierówność równoważnie (równoważnie, gdyż $ab+bc+ca>0$):
\begin{align*}
    a^2+b^2+c^2&<2ab+2bc+2ca,\\
    a^2+b^2+c^2+2ab-2bc-2ca&<4ab,\\
    (a+b)^2-2(a+b)c+c^2&<4ab,\\
    (a+b-c)^2&<4ab,\\
    |a+b-c|&<2\sqrt{ab}.
\end{align*}
Ostatnia nierówność oznacza, że
\[-2\sqrt{ab}<a+b-c<2\sqrt{ab},\]
czyli $a+b+2\sqrt{ab}>c$ oraz $a+b-2\sqrt{ab}<c$. Te warunki oznaczają odpowiednio, że
\[\sqrt{a}+\sqrt{b}>\sqrt{c}\quad\text{oraz}\quad \left|\sqrt{a}-\sqrt{b}\right|<\sqrt{c}.\]
Z~kolei drugi warunek jest równoważny układowi warunków 
\[-\sqrt{c}<\sqrt{a}-\sqrt{b}<\sqrt{c},\]
czyli $\sqrt{a}+\sqrt{c}>\sqrt{b}$ oraz $\sqrt{c}+\sqrt{b}>\sqrt{a}$. Uzyskaliśmy więc wszystkie trzy nierówności trójkąta dla liczb $\sqrt{a}$, $\sqrt{b}$, $\sqrt{c}$, z~których wynika teza zadania.

\paragraph{Rozwiązanie 4.} Podobnie jak w rozwiązaniu 3 przekształcamy założenie równoważnie do postaci
\[c^2-2(a+b)c+(a-b)^2<0\]
i~traktujemy ją jako kwadratową względem $c$ (z~parametrami $a$ oraz $b$). Wyróżnik trójmianu po lewej stronie jest równy
\[\Delta_c =4(a+b)^2-4(a-b)^2=16ab,\]
więc jego pierwiastkami są
\[c_1=\frac{2(a+b)+4\sqrt{ab}}{2}=\left(\sqrt a+\sqrt b\right)^2,\qquad c_2=\frac{2(a+b)-4\sqrt{ab}}{2}=\left(\sqrt{a}-\sqrt{b}\right)^2.\]
To zaś oznacza (skoro $|\sqrt{a}-\sqrt{b}|<\sqrt{a}+\sqrt{b}$), że nierówność jest spełniona wtedy i~tylko wtedy, gdy $c\in\left((\sqrt{a}-\sqrt{b})^2,(\sqrt{a}+\sqrt{b})^2\right)$, a~zatem
\[\sqrt{c}<\sqrt{a}+\sqrt{b}\quad\land\quad \sqrt{c}>|\sqrt{a}-\sqrt{b}|.\]
Drugi warunek jest z~kolei równoważny koniunkcji $\sqrt{c}>\sqrt{a}-\sqrt{b}$ oraz $\sqrt{c}>\sqrt{b}-\sqrt{a}$, czyli ostatecznie otrzymujemy
\[\sqrt{a}+\sqrt{b}>\sqrt{c}\quad\land\quad \sqrt{b}+\sqrt{c}>\sqrt{a}\quad\land\quad \sqrt{c}+\sqrt{a}>\sqrt{b},\]
a~to właśnie oznacza, że istnieje trójkąt o~bokach $\sqrt{a}$, $\sqrt{b}$, $\sqrt{c}$.

\paragraph{Komentarz} Zadanie można rozwiązać na wiele sposobów --- wspólną trudnością ich wszystkich są umiejętne przekształcenia algebraiczne (z~zastosowaniem wzorów skróconego mnożenia), umiejętna algebraizacja danego warunku (być może z~dodatkowym użyciem symetrii i~przyjęciem wygodnych założeń porządkujących $a$, $b$, $c$) oraz ostrożne wnioskowanie przy dowodzeniu nierówności. 

Warto zwrócić uwagę na to, że teza zadania ma postać implikacji: jeśli liczby $a$, $b$, $c$ spełniają warunek \eqref{2.23:teza}, to istnieje trójkąt. Nie kontrolując świadomie równoważności przekształceń w~rozwiązaniu 1 łatwo popełnić błąd logiczny polegający na wywnioskowaniu założenia z~tezy.

W rozwiązaniu 2. warto zwrócić uwagę na fakt, że dodatniość iloczynu czterech czynników nie jest równoważna dodatniości każdego z~nich (a:~parzystej liczbie ujemnych czynników i~braku zerowych). Potencjalny błąd polega więc na stwierdzeniu, że z~nierówności \eqref{2.23:iloczyn} wynika wprost teza zadania (nie wykluczając przypadku, gdy dwa lub cztery czynniki są ujemne).

Sztuczka ze zwinięciem w~rozwiązaniu 2 jest bardziej dostępna dla osób znających rozwiniętą postać wzoru Herona na pole trójkąta (lewa strona nierówności \eqref{2.23:iloczyn} to dokładnie $16P^2$, gdzie $P$ jest polem trójkąta o~bokach $x$, $y$, $z$.

Rozwiązanie numer 4 jest ładną ilustracją metody traktowania wyrażenia wielu zmiennych jako funkcji kwadratowej ze względu na jedną z~nich --- w~ten sposób uzyskujemy \emph{de facto} najbardziej szkolną metodę, sprowadzającą się do rozwiązania nierówności kwadratowej. Przy okazji można skontrolować, czy rozwiązujący nie popełni błędu przy ,,pierwiastkowaniu'' nierówności (czyli, czy nie zgubi modułu).

Do każdego rozwiązania pasuje inna seria wskazówek. Na przykład do rozwiązania numer 1 można rozważyć następujące:

\paragraph{Wskazówka 1.} W~jaki sposób można algebraicznie wyrazić tezę zadania?

\paragraph{Wskazówka 2.} Warunek jest symetryczny ze względu na $a$, $b$, $c$, więc można bez straty ogólności przyjąć np. $a\leq b\leq c$.

\paragraph{Wskazówka 3.} Spróbuj przekształcić tezę zadania tak, aby pozbyć się pierwiastków.

\paragraph{Uogólnienia i modyfikacje} Warto rozważyć zadanie pytania, jakie trójki (niekoniecznie dodatnich) liczb rzeczywistych $a$, $b$, $c$ spełniają równość
\[\frac{a^2+b^2+c^2}{ab+bc+ca}=2.\]
W~tym zadaniu dodatkowo trzeba zadbać o~wykluczenie zera w~mianowniku (i~ostrożnie kontrolować znaki --- np. argument z~deltą z~rozwiązania 4 pokazuje, że pośród $a$, $b$, $c$ nie może być dwóch o~ujemnym iloczynie).